\section{Форматы входных данных и отчётов}
\label{sec:formats}

\subsection{Начальный склад: CSV}
Файл, переданный через \code{--stock}, задаёт \emph{весь} набор начальных
партий. Лекарства, отсутствующие в файле, первоначально имеют нулевой остаток.
Каталог остаётся сгенерированным по параметрам эксперимента.

\begin{lstlisting}
medicine_id,count,expiration_day,unit_price
0,20,2,100
0,60,90,100
1,80,45,150
2,40,60
\end{lstlisting}

\begin{tabularx}{\textwidth}{@{}L{0.28\textwidth}Y@{}}
\toprule
Столбец & Семантика \\
\midrule
\code{medicine_id} & Целый индекс каталога от 0 до $K-1$.
Повторный индекс создаёт ещё одну партию. \\
\code{count} & Неотрицательное целое количество упаковок партии. \\
\code{expiration_day} & Неотрицательный последний день годности
в шкале эксперимента, а не календарная дата. \\
\code{unit_price} & Неотрицательная конечная закупочная цена упаковки.
Если столбец отсутствует, используется \code{wholesale_price} из каталога. \\
\bottomrule
\end{tabularx}

Допускаются пустые строки и комментарии от символа \texttt{\#} до конца строки.
Заголовок с \code{medicine_id} можно разместить до первой числовой партии.
В реализации запятые заменяются пробелами, после чего считываются числовые
поля; кавычки и сложные текстовые значения CSV не обрабатываются.
Ошибочная цена, отрицательное поле, лишний столбец или неизвестный
идентификатор приводят к исключению с указанием строки для части ошибок.
Пустой файл допустим и создаёт пустой начальный склад.

Для \code{expiration_day=0} партия может существовать при инициализации,
но списывается в начале дня 1. Цены из CSV используются для начальных партий;
новые поставки покупаются по оптовым ценам каталога.
Путь к CSV разрешается относительно текущей рабочей директории процесса.

\subsection{JSON: общая структура}
Отчёт создаётся функцией \code{model_detail::write_json_report(path)}
из \code{internal/simulation_cli.h}; консольный запуск вызывает её по
\code{--report} после завершения последнего дня. Корневой объект содержит:

\begin{tabularx}{\textwidth}{@{}L{0.28\textwidth}Y@{}}
\toprule
Ключ & Содержимое \\
\midrule
\code{parameters} & Объект параметров эксперимента. \\
\code{totals} & Итоговые денежные суммы, количества и состояния очередей. \\
\code{catalog} & Массив из $K$ сгенерированных лекарств. \\
\code{days} & Массив завершённых дней в порядке номеров. \\
\code{completed_orders} & Массив реально доставленных покупок. \\
\code{inventory} & Текущие партии склада на момент записи отчёта. \\
\code{pending_supplies} & Текущие ещё не поступившие заявки. \\
\bottomrule
\end{tabularx}

Строки записываются в UTF-8; кавычки, обратные косые черты и управляющие
символы экранируются. Суммы --- JSON-числа с точностью вывода 17 значащих цифр.
Индексы начинаются с нуля. Отдельного поля версии схемы и временной метки нет.
Следующая таблица описывает текущий формат, а не обещание его неизменности
при последующих изменениях программы.

\subsection{Объект параметров}
\code{parameters} включает:
\begin{lstlisting}
N, M, K, seed, markup_percent,
card_discount, bulk_discount, regular_discount, maximum_discount,
initial_count, regular_customers, initial_stock_file
\end{lstlisting}
\code{N}, \code{M}, \code{K} соответствуют полям C++ \code{days},
\code{couriers}, \code{medicine_count}. \code{bulk_discount} соответствует
\code{large_purchase_discount}. \code{initial_stock_file} --- строка пути,
пустая при автоматической генерации. \code{initial_count} сохраняется
в параметрах даже при использовании CSV и не заменяет фактическое
\code{totals.initial_packages}.

\subsection{Итоговые показатели}
{\small
\begin{longtable}{@{}L{0.37\textwidth}L{0.56\textwidth}@{}}
\toprule
Ключ в \code{totals} & Значение \\
\midrule
\endhead
\code{income} & Выручка всех доставленных покупок. \\
\code{purchase_expenses} & Стоимость начального запаса и прибывших поставок. \\
\code{write_off_losses} & Закупочная стоимость списанных упаковок. \\
\code{profit} & Прибыль по формулам раздела~\ref{sec:model}. \\
\code{available_stock_cost} & Закупочная стоимость доступных партий. \\
\code{reserved_stock_cost} & Закупочная стоимость уже выделенного,
но ещё не доставленного товара. \\
\code{delivered_stock_cost} & Закупочная стоимость всех доставленных упаковок. \\
\code{initial_packages} & Количество упаковок в начальных партиях до дня 1. \\
\code{supplied_packages} & Количество упаковок во всех прибывших поставках. \\
\code{written_off_packages} & Число упаковок, списанных по сроку годности. \\
\code{pending_deliveries} & Число подготовленных покупок на следующий день. \\
\code{pending_orders} & Число запросов, ещё ожидающих выделения товара. \\
\bottomrule
\end{longtable}
}

\subsection{Каталог}
Каждый элемент \code{catalog} содержит \code{medicine_id}, \code{dosage},
\code{type}, \code{group}, \code{wholesale_price}, \code{shelf_life},
\code{min_stock}. Поля \code{type} и \code{group} --- строки названий,
а не исходные индексы \code{type_med_id} и \code{group_med_id}.
Оптовая цена описывает сгенерированный каталог и может отличаться от
начальных цен конкретных партий CSV.

\subsection{Результат дня}
{\small
\begin{longtable}{@{}L{0.37\textwidth}L{0.56\textwidth}@{}}
\toprule
Ключ элемента \code{days} & Значение \\
\midrule
\endhead
\code{day} & Номер завершённого дня от 1 до $N$. \\
\code{received_orders} & Все новые запросы дня: плановые и обычные. \\
\code{planned_orders} & Плановая часть новых запросов. \\
\code{empty_orders} & Рассмотренные в этот день запросы с нулевой выдачей;
они могут поступить раньше. \\
\code{prepared_orders} & Успешно выделенные в этот день покупки на завтра. \\
\code{delivered_orders} & Покупки, доставленные сегодня после вчерашней подготовки. \\
\code{late_deliveries} & Из сегодняшних доставок: выполненные позже дня,
следующего за поступлением запроса. \\
\code{pending_orders} & Длина ещё не обработанной очереди в конце дня. \\
\code{arrived_supplies} & Число заявок, поступивших сегодня; это не число упаковок. \\
\code{created_supplies} & Число заявок, созданных сегодня. \\
\code{income}, \code{purchase_expenses}, \code{write_off_losses}
& Накопленные денежные суммы по конец дня включительно. \\
\code{courier_limits_satisfied} & Булево значение: у каждого курьера
от 7 до 15 реальных доставок. \\
\code{courier_load} & Массив длины $M$; элемент $c$ --- нагрузка курьера $c$. \\
\code{available_packages} & Массив длины $K$; элемент $i$ --- суммарный
остаток лекарства $i$ после выделения товара. \\
\bottomrule
\end{longtable}
}

Показатель выручки отдельного дня вычисляется разностью соседних
\code{income}. Для закупочных расходов дня 1 вычитается стоимость начального
запаса, если требуется отделить закупки до эксперимента от поступлений первого
дня. В \code{available_packages} товар на завтрашние доставки уже отсутствует.

\subsection{Доставленная покупка}
Каждый элемент \code{completed_orders} содержит идентификатор
\code{order_id}, дни \code{ordered_day}, \code{prepared_day},
\code{delivered_day}, \code{courier_index}, \code{price},
\code{acquisition_cost}, объект \code{customer} и массивы
\code{requested} и \code{supplied}.

В JSON-объект покупателя входят \code{fio}, \code{phone}, \code{address},
\code{has_card}, \code{is_regular}, \code{card_number}.
Расписание регулярных покупок в JSON не сериализуется, хотя присутствует
в копии \code{Customer} результата C++.
Позиция каждого массива имеет форму:
\begin{lstlisting}
{"medicine_id": 0, "count": 2}
\end{lstlisting}

Пример одного элемента ниже иллюстрирует схему. Значения условны:
\begin{lstlisting}
{
  "order_id": 1,
  "ordered_day": 1,
  "prepared_day": 1,
  "delivered_day": 2,
  "courier_index": 0,
  "price": 125,
  "acquisition_cost": 100,
  "customer": {
    "fio": "Example customer",
    "phone": "+79000000000",
    "address": "Example address",
    "has_card": false,
    "is_regular": false,
    "card_number": 0
  },
  "requested": [{"medicine_id": 0, "count": 2}],
  "supplied": [{"medicine_id": 0, "count": 1}]
}
\end{lstlisting}
Такая запись отражает частичную выдачу одной упаковки при запросе двух.
Отсутствующее количество не создаёт отдельного обязательства доставки.

\subsection{Партии и ожидающие поставки}
Каждый элемент \code{inventory} содержит \code{medicine_id}, \code{count},
\code{expiration_day}, \code{unit_price}. Здесь \code{unit_price} ---
\emph{текущая продажная база} \code{batch_price}: она уже может быть уменьшена
уценкой. В CSV одноимённый столбец задавал первоначальную закупочную цену.
Отдельные закупочные цены остатка в этот массив не экспортируются.
После выделения товара в конце дня в векторе могут оставаться партии
с нулевым количеством до следующего вызова обработки склада.

Каждый элемент \code{pending_supplies} содержит \code{medicine_id},
\code{count} и \code{delivery_timer}. Таймер обозначает оставшиеся дни
до исполнения заявки, а не её исходную случайную задержку.
День заявки, ожидаемая дата и полная история поступлений в JSON не включены.

\subsection{Объём экспортируемых сведений}
JSON содержит подробную историю \emph{доставленных} покупок и дневные
суммарные остатки. Подготовленные покупки представлены количеством и
общей стоимостью, очередь необработанных запросов --- количеством.
Их поштучное содержимое, история партий по каждому дню и причины каждой
частичной выдачи не экспортируются. Для такого анализа потребовалось бы
расширение структур результатов и сериализатора.

Для внешней проверки доступны равенства:
\begin{align}
  \sum_{o\in\mathrm{completed}} o.\mathrm{price}
    &= \mathrm{totals.income},\\
  \sum_{c=0}^{M-1}\mathrm{load}_{d,c}
    &= \mathrm{delivered\_orders}_d,\\
  O_{\mathrm{received}}&=O_{\mathrm{prepared}}+O_{\mathrm{empty}}
                        +O_{\mathrm{pending}}.
\end{align}
В последнем равенстве первые три величины --- суммы соответствующих
счётчиков за весь период, а $O_{\mathrm{pending}}$ --- итоговая длина очереди.
Денежные сравнения следует выполнять с допуском для \code{double}.
Число подготовленных покупок за весь период равно числу выполненных
плюс \code{pending_deliveries}.
